% Einheit 2 von 4 – Satz vom Nullprodukt (bank/quadratische-gleichungen/e2.jsonl, zusammenbau v0.1)
%% TODO Zweigzeile Teil 1: was hier gelernt wird (Satz in Schülersprache aus dem Einheitstitel) – Einheit 2
\einheitenkopf[e2]{Einheit 2 von 4 · Satz vom Nullprodukt}
\zweigzeile{ab Kl. 10 (am Gymnasium ab Kl. 8) · P10}
\setcounter{aufgabe}{13}

%% TODO Titel: Ich-kann-Satz fehlt in der Bank (Platzhalter: Kettenname) – Nr. 14
%% TODO Anweisung: ein Satz über der ersten Teilaufgabe fehlt in der Bank – Nr. 14
\begin{aufgabe}{Nullprodukt}
\begin{teile}
\teil $(x - 2) \cdot (x + 7) = 0$ – gilt der Satz vom Nullprodukt?\\ \kreuz{gilt} \\ \kreuz{gilt nicht}
\end{teile}
%% TODO Darstellung für schwach fehlt in der Bank (quadratische-gleichungen-e2-k1-s1-v1; Abschnitt „Für schwache Schüler“)
\swz{$\text{(x - 2) \cdot (x - 5) = 0}$}{}
%% TODO Darstellung für schwach fehlt in der Bank (quadratische-gleichungen-e2-k1-s1-v2; Abschnitt „Für schwache Schüler“)
\swz{$\text{(x - 1) \cdot (x - 7) = 0}$}{}
%% TODO Darstellung für schwach fehlt in der Bank (quadratische-gleichungen-e2-k1-s1-v3; Abschnitt „Für schwache Schüler“)
\swz{$\text{(x - 3) \cdot (x - 10) = 0}$}{}
%% TODO Darstellung für schwach fehlt in der Bank (quadratische-gleichungen-e2-k1-s1-v4; Abschnitt „Für schwache Schüler“)
\swz{$\text{(x - 6) \cdot (x - 4) = 0}$}{}
%% TODO Darstellung für schwach fehlt in der Bank (quadratische-gleichungen-e2-k1-s1-v5; Abschnitt „Für schwache Schüler“)
\swz{$\text{(x - 8) \cdot (x - 3) = 0}$}{}
\swfrage{Was bleibt gleich, was ändert sich?}
%% TODO Darstellung für schwach fehlt in der Bank (quadratische-gleichungen-e2-k1-s2-v1; Abschnitt „Für schwache Schüler“)
\swz{$\text{(x + 4) \cdot (x - 1) = 0}$}{}
%% TODO Darstellung für schwach fehlt in der Bank (quadratische-gleichungen-e2-k1-s3-v1; Abschnitt „Für schwache Schüler“)
\swz{$\text{x \cdot (x + 6) = 0}$}{}
%% TODO Darstellung für schwach fehlt in der Bank (quadratische-gleichungen-e2-k1-s4-v1; Abschnitt „Für schwache Schüler“)
\swz{$\text{(x - 7) \cdot (x - 7) = 0}$}{}
%% TODO Darstellung für schwach fehlt in der Bank (quadratische-gleichungen-e2-k1-s5-v1; Abschnitt „Für schwache Schüler“)
\swz{Ist $x = -5$ Lösung von $(x + 5) \cdot (x - 1) = 0$?}{}
%% TODO Darstellung für schwach fehlt in der Bank (quadratische-gleichungen-e2-k1-s6-v1; Abschnitt „Für schwache Schüler“)
\swz{$\text{x^2 + 9x = 0}$}{}
%% TODO Darstellung für schwach fehlt in der Bank (quadratische-gleichungen-e2-k1-s7-v1; Abschnitt „Für schwache Schüler“)
\swz{$\text{x^2 - 3x = 0}$}{}
%% TODO Darstellung für schwach fehlt in der Bank (quadratische-gleichungen-e2-k1-s8-v1; Abschnitt „Für schwache Schüler“)
\swz{$\text{2x^2 - 6x = 0}$}{}
%% TODO Darstellung für schwach fehlt in der Bank (quadratische-gleichungen-e2-k1-s9-v1; Abschnitt „Für schwache Schüler“)
\swz{$x \cdot (x + 3) = 10$ – ausmultiplizieren und ordnen?}{}
\begin{teile}
\teil Welcher Wert erfüllt $x \cdot (x + 9) = -14$? \hfill (P10 2025 OS)\\ \kreuz{$x = 2$} \\ \kreuz{$x = 7$} \\ \kreuz{$x = -7$} \\ \kreuz{$x = -14$}
\end{teile}
\end{aufgabe}

%% TODO Titel: Ich-kann-Satz fehlt in der Bank (Platzhalter: Kettenname) – Nr. 15
%% TODO Anweisung: ein Satz über der ersten Teilaufgabe fehlt in der Bank – Nr. 15
\begin{aufgabe}{Nullprodukt – Fehler finden, Begründen}
\swz[3]{Wo ist der Fehler? \rechnung{x^2 &= 6x &&\mid : x \\ x &= 6}}{}
\swfrage{Begründe: Ist bei $(x - 3) \cdot (x + 5)$ einer der beiden Faktoren null, ist das ganze Produkt null.}
\end{aufgabe}

\uebersichtskasten{Nullprodukt: Ein Produkt ist null, sobald ein Faktor null ist. Steht rechts null, jeden Faktor einzeln null setzen. \\ x² − 7x = 0 → x·(x − 7) = 0 → x₁ = 0, x₂ = 7 \quad (x + 3)·(x − 9) = 0 → x₁ = −3, x₂ = 9 \\ Ohne Absolutglied immer x ausklammern – nie durch x teilen, sonst geht die Lösung null verloren. \\ Steht rechts nicht null, gilt der Satz nicht: ausmultiplizieren und ordnen. \quad x·(x + 8) = 20 → x² + 8x − 20 = 0 → Formel (Einheit 3)}
